\section{Conclusion}

The 2030 Sun Belt redistricting challenge requires two structural changes to the
bisect pipeline:
\begin{enumerate}
  \item For TX and FL: switch from ApportionRegions (prime-factor) to
    GeoSection (\texttt{--structure ratio-optimal}). GeoSection handles any $k$
    and produces better geographic alignment for the irregular geometry of
    fast-growing coastal and urban states.
  \item For seat-loss states (NY, CA, IL, OH, MI, PA, others): rebuild the
    bisection tree from scratch for the new $k$, without any continuity
    constraint from the 2020 plan. This is a feature, not a limitation: the
    algorithmically optimal 2030 plan is better than an incumbency-protecting
    merger.
\end{enumerate}

The key algorithmic insight: the binary fallback for prime $k$ in ApportionRegions
(lines 11--16 of prime.rs) means no state will ever require a direct $k$-way
METIS partition for prime $k > 3$. This design choice makes the pipeline robust
to any projected seat count, including TX $k=41$.

All bisection tree designs for the 2030 seat allocations should be pre-built
and validated by 2029, well before the April 2031 Census data release.
